% Generated draft: requires mathematical and source review.
\documentclass[11pt]{article}
\usepackage[margin=25mm]{geometry}
\usepackage{fontspec}
\setmainfont{Latin Modern Roman}
\usepackage{amsmath,amssymb,mathtools,unicode-math}
\setmathfont{Latin Modern Math}
\usepackage{enumitem,needspace,xurl,hyperref}
\setlength{\parindent}{0pt}
\setlength{\parskip}{5pt}
\setlength{\emergencystretch}{3em}
\newcommand{\N}{\mathbb{N}}
\newcommand{\Z}{\mathbb{Z}}
\newcommand{\Q}{\mathbb{Q}}
\newcommand{\R}{\mathbb{R}}
\newcommand{\C}{\mathbb{C}}
\newcommand{\F}{\mathbb{F}}
\newcommand{\A}{\mathbb{A}}
\newcommand{\PP}{\mathbb{P}}
\newcommand{\E}{\mathbb{E}}
\newcommand{\eps}{\varepsilon}
\newcommand{\dd}{\,\mathrm{d}}
\begin{document}
\section*{UnsolvedMath ID: 20001280}
\textbf{Source problem code:} catalogue entry
\subsection{Definitions and mathematical statement}
The verification identifies this record as a concatenation of Problems 13--16, page 2 of the 2010 AIM workshop list, despite the extracted opening number ``3.'' Let \(K\subset\R^n\) be an origin-symmetric convex body, meaning compact, convex, with nonempty interior, and let
\[
K^\circ=\{y\in\R^n:x\cdot y\le1\text{ for every }x\in K\}
\]
be its polar (the source writes \(K^*\)).

Problem 13 (T. Tao) asks whether the polar has the largest lattice covering radius when \(K\) is a cube. Its literal displayed radius deletes the lattice centers lying in \(K\):
\[
r_{\rm del}(K)=\inf\{r>0:\ \bigcup_{v\in\Z^n\setminus K}(v+rK^\circ)=\R^n\}.
\]
This differs from the standard covering radius
\[
\mu(K^\circ,\Z^n)=\inf\{r>0:\ \bigcup_{v\in\Z^n}(v+rK^\circ)=\R^n\}.
\]
The source supplies no size, lattice-position or affine normalization. The aid reports that the proposed unrestricted maximum is untenable under either reading. For each fixed \(K\), deletion of the finitely many centers \(K\cap\Z^n\) still gives \(r_{\rm del}(K)<\infty\). If \(aB_2^n\subseteq K\), then, for \(\lambda>0\),
\[
r_{\rm del}(\lambda K)\ge\lambda^2a^2,
\qquad
\mu((\lambda K)^\circ,\Z^n)=\lambda\mu(K^\circ,\Z^n).
\]
In the deleted-center expression the forbidden centers change with \(\lambda K\); they are not held fixed during dilation.

Even a volume normalization leaves the quantities unbounded in dimensions \(n\ge2\). For \(R\ge1\), set
\[
K_R=[-R,R]\times[-R^{-1},R^{-1}]\times[-1,1]^{n-2}.
\]
These determinant-one linear images of the cube all have volume \(2^n\). The aid gives
\[
\mu(K_R^\circ,\Z^n)=\frac{R+R^{-1}+n-2}{2},\qquad r_{\rm del}(K_R)\ge R/2.
\]
Thus a meaningful extremal version needs an additional intended restriction, which is not supplied; this draft does not invent one.

For the appended questions, the intersection-body operator is given, on origin-symmetric star bodies, by
\[
\rho_{IK}(u)=\operatorname{vol}_{n-1}(K\cap u^\perp),\qquad u\in S^{n-1},
\]
where \(\rho_K(u)=\max\{t\ge0:tu\in K\}\). Problem 14 asks, for \(n>3\), whether \(IK=K\) implies \(K=cB_2^n\) for some \(c>0\). The equality is literal, so a ball that is fixed also has the corresponding radius normalization; the conclusion only asserts spherical shape. The source notes that the analogous spherical-shape conclusion fails for projection bodies because a cube's projection body is a dilate of a cube. The usual projection-body convention is \(h_{\Pi K}(u)=\operatorname{vol}_{n-1}(K|u^\perp)\). It calls the intersection question a special case of [G, Problem 8.7], refers to [G, Note 8.6], and cites [G, Problem 4.5] and [G, Note 4.6] for the more general projection-body question and Weil's solution for polytopes. The aid reports that Problem 14 was solved affirmatively in 2025, but supplies no author or publication citation for that status claim.

Problem 15 asks, for \(n>3\), whether the iterates \(I^mK\) converge to a Euclidean ball in Banach--Mazur distance as \(m\to\infty\), equivalently
\[
\lim_{m\to\infty}d_{\rm BM}(I^mK,B_2^n)=1.
\]
For symmetric convex bodies, \(d_{\rm BM}(A,B)\) is the infimum of \(\lambda\ge1\) for which \(A\subseteq TB\subseteq\lambda A\) for some invertible linear map \(T\); it ignores scalar dilation. This is not literal convergence of unnormalized bodies. The domain of \(K\) in the appended questions is not restated; retaining the opening symmetric-convex-body hypothesis is an explicit interpretation. Any broader star-body version would need its distance convention specified.

Problem 16 asks, for \(n>5\), to construct a polytope \(K\) that is an intersection body and is not the polar of a zonotope. A zonotope is a finite Minkowski sum of line segments. For an origin-symmetric \(K\), the exclusion can be written as \(K^\circ\) not being a zonotope. The source does not say whether ``intersection body'' means exactly \(IL\) for a star body \(L\), or membership in a closure/generalized class; that distinction must be checked before choosing a construction.

The aid describes Problems 15 and 16 as globally open in the literature it checked. This is a supplied literature qualification, not an independent current assessment. The registry labels the combined entry partially solved and reconstructed unverified; attempt 1 is a partial result, with the problem partially solved at the run, internal rigor proved and novelty confidence low. The research aid is machine generated, not peer reviewed, and requires independent verification before citation.
\subsection{Short English statement}
Four questions were accidentally joined into one record. The covering-radius question has an unspecified normalization and is unbounded even at fixed volume under both its literal deleted-center formula and the usual formula. The appended questions concern fixed points and iterates of the intersection-body operator and polytopal intersection bodies; the aid reports a 2025 solution only for the fixed-point question.
\subsection{Sources}
\begin{itemize}
\item Local UnsolvedMath registry, permanent ID \(20001280\), source code \texttt{catalogue entry}; \texttt{aim-convex-geometry-notes.json}, \texttt{notes[11]}.
\item AIM, \emph{Mahler's conjecture and duality in convex geometry} (2010), page 2, recovered Problems 13--16; Problem 13 attributed to T. Tao: \url{https://aimath.org/WWN/mahlerduality/mahlerduality.pdf}. The recovery is reported by the supplied verification, not independently performed here.
\item Supplied references [G, Problem 8.7], [G, Note 8.6], [G, Problem 4.5], [G, Note 4.6], including the attribution to Weil for polytopes. The record does not identify [G] bibliographically; no author, book title or edition is inferred.
\item Supplied research attempt 1 and registry update August 14, 2026. No authoritative citation is supplied for the reported affirmative solution of Problem 14 in 2025 or for the claimed checked-literature status of Problems 15--16. The original literature field is empty.
\end{itemize}
\end{document}
